\section{Foundations of Agentic Visual Generation}
\label{sec:foundations}

\subsection{Visual Generation Foundations}
\label{sec:visual-foundations}

Visual generation in this survey covers the construction and editing of images, videos, three-dimensional scenes, parametric CAD models, scientific graphics, structured documents, interactive interfaces, and related visual products. Four foundations connect these operations to agentic control: the generator, the artifact representation, the conditioning inputs, and the execution interfaces. Together, they determine how an artifact is produced, what task-relevant information persists, what information an operation can use, and what actions the system can execute.

\subsubsection{Visual Generators}

For a stochastic base generator, the output artifact $x$ can be represented as a sample from a conditional distribution $p_{\phi}(x \mid c)$. The condition $c$ may encode textual instructions, visual or spatial references, structured task information, or state carried from earlier operations. Different generator families provide different ways to represent this conditional distribution and produce an artifact from the available task information. The following progression outlines the principal modeling families that underpin current visual generation and editing systems \cite{he2024llms,yang2025textimage,wang2025imageediting,yin2025spatiotemporal,wu2026visualgeneration}.

Early generative models introduced foundational approaches for representing and learning data distributions. Variational autoencoders used a latent-variable formulation with approximate inference, and generative adversarial networks learned a generator through competition with a discriminator \cite{kingma2013autoencoding,goodfellow2014gan}. Normalizing flows used invertible transformations with tractable density evaluation, while energy-based models represented compatibility through an energy function \cite{he2024llms,yang2025textimage}. These approaches established alternative parameterizations for later conditional generators.

Autoregressive generation made the ordering of visual content explicit. VQ-VAE and VQGAN showed how images can be represented as discrete visual tokens, providing a basis for token-level generation and editing \cite{oord2017vqvae,esser2021vqgan}. PixelRNN predicted pixels sequentially, and DALL-E demonstrated large-scale text-conditioned generation with discrete visual representations \cite{oord2016pixelrnn,ramesh2021dalle}. MaskGIT provided a related masked-token formulation: groups of tokens are filled in parallel and refined over multiple iterations \cite{chang2022maskgit}. Visual Autoregressive Modeling (VAR) arranged prediction from coarse to fine scales, using next-scale prediction to reduce the effective sequence length \cite{tian2024var}. These models organize visual generation as sequential or partially parallel prediction over pixels or visual tokens.

Recent autoregressive work revisits this route at larger scale. LlamaGen applies next-token prediction to visual tokens, while MAR removes the requirement for a discrete vector-quantized tokenizer by modeling continuous token distributions with a diffusion loss \cite{llamagen2024,mar2024}. A related line unifies multimodal understanding and generation in a single Transformer: Chameleon and Emu3 model mixed visual and textual sequences, Show-o combines autoregressive and discrete-diffusion objectives, and Janus uses separate visual encoding paths for understanding and generation \cite{chameleon2024,emu32024,showo2024,janus2024}.

Diffusion and score-based methods formed a closely related line of development. Denoising Diffusion Probabilistic Models (DDPM) described generation through the reversal of a discrete noising process. Score-based models learned the score of perturbed data distributions, and continuous-time stochastic differential equation formulations connect the two descriptions through a probability-flow ordinary differential equation \cite{ho2020ddpm,song2021scorebased}. Large-scale text-conditioned systems such as Imagen and latent diffusion moved high-resolution synthesis toward practical language-guided generation \cite{saharia2022imagen,rombach2022ldm}. Diffusion Transformers (DiT) introduced Transformer-based denoisers operating on latent patches \cite{peebles2023dit}. Flow matching and rectified flow learned continuous transport paths, while consistency models targeted one-step or few-step generation \cite{lipman2022flowmatching,liu2022rectified,song2023consistency}. Efficiency-oriented variants such as Sana illustrate later efforts to reduce the cost of high-resolution diffusion Transformers \cite{sana2024}. Conditional generation and editing systems build on these families by using text, references, regions, and other task inputs \cite{yang2025textimage,wang2025imageediting}.

Recent domain-specific generators extend these backends beyond static images. Video Diffusion Models established a diffusion-based route for temporal synthesis, while CogVideoX and subsequent open technical systems such as HunyuanVideo and Wan illustrate recent large-scale video generation \cite{ho2022videodiffusion,cogvideox2024,hunyuanvideo2024,wan2025}. Three-dimensional generation combines learned scene representations, image-conditioned reconstruction, and distillation from 2D priors: NeRF provided a neural scene representation, DreamFusion used a 2D diffusion prior for text-to-3D synthesis, and LRM and TripoSR demonstrated fast feed-forward reconstruction from a single image \cite{mildenhall2020nerf,poole2022dreamfusion,lrm2023,triposr2024}. Structured visual generation follows another route. DeepCAD represents CAD designs through editable construction operations, while later systems connect visual or language intent to CAD programs and geometric constraints \cite{wu2021deepcad,ocker2025ideatocad}. Related work generates layouts, data-driven charts, documents, or interface code, with LayoutDM, Data2Vis, and pix2code providing representative examples of structured output generation \cite{inoue2023layoutdm,dibia2018data2vis,beltramelli2017pix2code}. These backends provide the generator layer on which the artifact, conditioning, and execution foundations operate.

\subsubsection{Visual Artifacts and Task Evidence}

A visual artifact is the object that a system creates, edits, or delivers. Common forms include rendered images and videos, structured scene or geometric representations, data-grounded graphics, documents, interactive interfaces, and related artifacts. Depending on its representation, an artifact may contain rendered appearance together with temporal relations, geometry, parameters, source data, hierarchy, or interaction state. These representations can preserve task-relevant properties across operations and make them available for later inspection or modification \cite{yin2025spatiotemporal,mildenhall2020nerf,wu2021deepcad,dibia2018data2vis,beltramelli2017pix2code}.

The artifact representation, together with the task requirements, helps determine what constitutes a valid result. Video tasks may require temporal preservation of identities, events, motion, camera changes, and other sequence-level properties. Three-dimensional and CAD representations can expose geometry, topology, parameters, and construction constraints. Data-driven graphics may retain links to source data, structured documents may retain layout or hierarchy, and interfaces may retain executable behavior. Other visual products impose their own representation-specific requirements and sources of evidence. Correctness therefore depends on evidence matched to the representation and the task, such as temporal checks, geometric or constraint checks, data consistency, layout inspection, or functional execution. The representation and the available operations together determine which properties can be inspected, preserved, or modified later.

\subsubsection{Conditioning Inputs}

A conditioning input is information supplied before a particular generation or editing operation. It helps specify what that operation should create, preserve, or change. The main forms are textual instructions, visual or spatial references, and structured task information. Text provides semantic intent and requested content; visual or spatial references provide appearance, identity, location, or layout cues; structured information provides explicit data values, geometric constraints, programmatic specifications, or other task requirements \cite{yang2025textimage,wang2025imageediting,yin2025spatiotemporal}.

The role of a condition depends on the operation that receives it. An initial operation may use a user request and reference material, while a later operation may use the current artifact, a selected region, an intermediate specification, or a verified state. These inputs make some task variables explicit while leaving other requirements to be inferred, clarified, or checked after execution. They also affect the granularity of control: a global description may guide the whole artifact, while a region or object may support a localized change.

\subsubsection{Execution Interfaces}

An execution interface is the callable boundary through which a controller invokes an operation on an artifact or task environment, or obtains task-relevant execution evidence. It defines the callable operations, accepted inputs, returned outputs, and execution status or error information made available to the controller. For a structured artifact, an operation may alter a representation or render it; for an interactive interface or other executable artifact, it may run code, tests, or interaction steps. Interfaces differ in operation granularity and feedback channels. They connect a controller decision to an artifact or environment update, or to task-relevant execution evidence \cite{he2024llms,yin2025spatiotemporal}.

Conditions specify the information and requirements for an operation, whereas execution interfaces specify how that operation can be carried out. The available operations and their granularity determine whether a controller can make a global change, edit a local region or object, modify a parameter, call a verifier, or invoke a recovery operation. Their outputs determine what the controller can observe after execution, including a new artifact, a rendering, a structured result, an error, or a test outcome. Together, these properties define the action space and the execution evidence available at each point in a visual creation process.


\subsection{Definition of Agentic Visual Generation}
\label{sec:agentic-visual-generation}

\subsubsection{Behavioral Definition}

Agentic visual generation is identified by how a system controls a visual creation trajectory. It operates through a closed loop when runtime evidence from the evolving artifact, task environment, interaction history, or other task-relevant sources informs subsequent creation decisions. Depending on the task, these decisions may concern clarification, planning, retrieval, model or tool selection, generation, editing, verification, revision, recovery, termination, and related operations \cite{xie2024large,durante2024agentai,jiang2024rationality,schneider2025generative}.

At a practical level, AVG is goal-driven closed-loop visual creation: the system uses evidence available while the task is running to determine later actions. The formal definition is:

\begin{quote}
\emph{Agentic Visual Generation (AVG) denotes goal-driven closed-loop visual creation in which an identifiable decision process maintains accessible, task-relevant state for an evolving visual artifact, its task environment, and interaction history, and uses runtime evidence to select or parameterize subsequent creation actions.}
\end{quote}

The definition is applied in this survey by examining a system's creation trajectory. The task concerns the creation, editing, or construction of a visual artifact or editable visual environment. An identifiable decision process selects or parameterizes subsequent creation actions and retains task-relevant state in a form accessible at later decision points. Runtime evidence available during task execution, including evidence from intermediate artifacts, execution outcomes, external information, user feedback, and other task-relevant sources, is shown to alter a later action selection or parameterization. The decision process may be implemented by a single model, a learned policy, a search procedure, or a hierarchy of specialized agents. This survey applies the same trajectory-level criterion across these implementations.

\subsubsection{Formalization of Closed-Loop Visual Creation}

The four foundations above specify what can be produced, what representation is retained, what information is supplied to an operation, and how that operation is invoked. We abstract a visual-creation trajectory using these same elements. Let $g$ denote the task goal and its requirements, $x_t$ the current artifact representation, $s_t$ the relevant state of the task environment, $h_t$ the accessible interaction history up to step $t$, and $\mathcal{I}_t$ the execution interfaces available at step $t$. The artifact representation $x_t$ may be a rendered image or video, a structured scene or geometric model, a document or layout, executable interface code, or another representation appropriate to the task. The environment state $s_t$ includes information needed to continue execution, such as application state, source data, files, or execution results. The history $h_t$ records prior actions, observations, outcomes, and user interactions.

At each decision point, the controller receives runtime evidence about the current trajectory. We denote this task-dependent observation by $o_t$:

\begin{equation}
  o_t = \mathsf{O}(g,x_t,s_t,h_t).
\end{equation}

Here $o_t$ is a general evidence variable. It may contain a visual inspection, a structured property, an execution result, a constraint check, external information, or user feedback, depending on the artifact and task. The controller uses this evidence to select an action. When the action invokes an execution interface, it can be written as $a_t=(i_t,c_t)$, where $i_t\in\mathcal{I}_t$ is the selected interface and $c_t$ is the condition and operation input supplied to it. The action can also be a control decision such as revising the plan, recovering from an error, or stopping:

\begin{equation}
  a_t \sim \pi_{\theta}\!\left(\cdot \mid g,o_t,h_t,\mathcal{I}_t\right).
\end{equation}

The interface set $\mathcal{I}_t$ is the callable boundary described above: it determines which operations can be invoked, which conditions they accept, and which artifact updates, execution outcomes, or status information they return. When $a_t=(i_t,c_t)$ invokes a stochastic generator, its output can be written as $x_{t+1}\sim p_{\phi}(\cdot\mid c_t)$ for a new artifact. For editing and structured construction, the same interface updates an existing representation. Both cases are represented by the transition:

\begin{equation}
  (x_{t+1},s_{t+1},h_{t+1})
  = \mathsf{U}\!\left(g,x_t,s_t,h_t,a_t;\mathcal{I}_t\right).
\end{equation}

The updated artifact and environment then produce new runtime evidence, which is available to the next decision. When an evidence-producing check is invoked through $\mathcal{I}_t$, its result is recorded in $s_{t+1}$ or $h_{t+1}$ and becomes part of $o_{t+1}$ through $\mathsf{O}$. Its output can be visual, symbolic, geometric, data-based, executable, or interaction-based. Verification is thus tied to the representation and task requirements established in the foundation section.

The trajectory is closed-loop when runtime evidence changes later action selection. Formally, holding the task, prior trajectory, and available interfaces fixed, there must be a decision point at which different runtime evidence leads to different subsequent action distributions:

\begin{equation}
  \exists\,t,o\ne o':\quad
  \pi_{\theta}\!\left(\cdot\mid g,o,h_t,\mathcal{I}_t\right)
  \ne
  \pi_{\theta}\!\left(\cdot\mid g,o',h_t,\mathcal{I}_t\right).
\end{equation}

\subsection{Boundary of Agentic Visual Generation}
\label{sec:agentic-visual-generation-boundary}

The boundary distinguishes related visual-creation systems by their control paths. A one-shot generator maps initial conditions to an artifact, while a fixed workflow carries artifacts through predetermined stages. An agent-augmented workflow may add tool, model, or role choices within that process. The control path becomes agentic when information obtained during execution changes a later creation decision. The following representative types describe this progression and their relation to AVG.

\begin{table}[t]
  \caption{Representative visual-creation system types and their relation to AVG.}
  \label{tab:boundary}
  \small
  \begin{tabularx}{\columnwidth}{@{}p{0.25\columnwidth}p{0.43\columnwidth}p{0.19\columnwidth}@{}}
    \toprule
    System type & Control path & Relation to AVG \\
    \midrule
    One-shot generator & Single condition-to-artifact mapping & Not AVG; foundation and baseline \\
    Fixed workflow & Predetermined stages, tools, and parameters; earlier outputs may pass to later stages without changing the schedule & Not AVG; pipeline baseline \\
    Agent-augmented workflow & Limited tool, model, or role choices within a prescribed process & Not necessarily AVG; AVG only when runtime evidence changes a later action \\
    Closed-loop visual agent & Runtime observations change generation, editing, verification, recovery, stopping, or related creation decisions & AVG \\
    Continually improving agent & A closed-loop visual agent additionally changes reusable memory, skills, verifiers, policies, or models across tasks & AVG extension; cross-task improvement is not required for AVG \\
    \bottomrule
  \end{tabularx}
\end{table}

The next subsection organizes the extent of state-dependent control and cross-task change.

\subsection{Autonomy Levels}
\label{sec:agentic-visual-generation-autonomy}

AVG systems can support different spans of state-dependent control. We use five cumulative levels to record the strongest behavior supported by the available evidence, from a fixed baseline to cross-task improvement.

\begin{table}[t]
  \caption{L1--L5 autonomy levels used throughout the survey.}
  \label{tab:autonomy}
  \small
  \begin{tabularx}{\columnwidth}{@{}llX@{}}
    \toprule
    Level & Name & Operational criterion \\
    \midrule
    L1 & Fixed mapping or pipeline & The output and control path are predetermined. \\
    L2 & Tool or role assistance & The system selects tools, models, or roles without a demonstrated artifact-dependent feedback loop. \\
    L3 & Feedback adaptation & Intermediate evidence causes one or more subsequent revisions. \\
    L4 & Long-horizon autonomy & Persistent state supports dynamic replanning, failure recovery, and adaptive stopping. \\
    L5 & Continual self-improvement & Experience changes reusable memory, skills, verifiers, policy, or model behavior across tasks. \\
    \bottomrule
  \end{tabularx}
\end{table}

L1 is the fixed, non-adaptive baseline. L2 exposes an action space through tool selection, routing, or role coordination. L3 requires an observable dependency between intermediate evidence and a later revision. L4 extends this dependency across multiple operations through persistent state, replanning, recovery, and stopping. L5 additionally requires transfer: experience from earlier tasks changes later behavior under held-out evaluation.
